% =====================================================================================
% sections/S1_problems.tex -- Supplementary Section S1: the test problems (statements of the
% existence and uniqueness results, evaluation sets, further checks of the exact solutions).
% All `% src:` paths are relative to the root of the repository.
% =====================================================================================
\section{Test problems: exact solutions, uniqueness and evaluation sets}\label{S1:sec}

This section completes Section~\ref{s2_problems:sec}. A classical solution is one that is
twice continuously differentiable on the closed domain (the closed space--time cylinder for
\PH, \PWone\ and \PWtwo). The proofs of the two propositions are in
Section~\ref{sA_proofs:sec}.

\begin{proposition}[Exact solutions]\label{s2_problems:prop:exact}\leavevmode
\StmtExact
\end{proposition}

\begin{proposition}[Uniqueness]\label{s2_problems:prop:unique}\leavevmode
\StmtUnique
\end{proposition}

Propositions~\ref{s2_problems:prop:exact} and~\ref{s2_problems:prop:unique} together say that
each problem has a unique solution in the stated class and that it is the function $\uex$
against which we measure; for \PLthree\ the class is the one in which the datum is attained
in mean square.

\paragraph{Evaluation sets.}
% src: data/s2_problems_evalsets.json (script code/s2_problems_evalsets.py):
%      heat1d, wave1d, laplace2d: 4 of 40401 evaluation nodes coincide with boundary/initial
%      training points; wave2d: 1241 of 137781 (fraction 0.0090); laplace3d: 0 of 64000;
%      cell-centred collocation grid (grid-c) coinciding with evaluation points: 0 for
%      heat1d, wave1d, wave2d, laplace2d; 512 of 64000 (fraction 0.008) for laplace3d.
% src: research_highdim/scripts/hd_core.py (fixed_sets: separate generators for the test sets)
The evaluation sets are dense grids, or random points drawn from a generator of their own.
They are separate from the training points with two small exceptions. Grid nodes on the
boundary or the initial line can coincide with boundary or initial training points: the four
corner nodes for \PH, \PWone\ and \PLtwo\ (of $40\,401$), and $1241$ of the $137\,781$ nodes
for \PWtwo\ ($0.9$ per cent). And for \PLthree\ the $512$ collocation points of the
cell-centred strategy \strat{grid-c} are among the $64\,000$ evaluation points ($0.8$ per
cent); at those points the network is trained on the residual, not on values of $\uex$.

\paragraph{Further checks of the exact solutions.}
The formulas were also checked by means that do not rely on the proofs.
% src: data/s2_problems_symbolic.txt (script code/s2_problems_symbolic.py; all 55 expressions 0)
(i)~Symbolic differentiation returns a zero residual and a zero mismatch in every initial
and boundary condition for the seven closed-form solutions (for \Pone\ and \Ptwo\ at
$d=2,3,5,10,20$); for \PLthree\ it confirms that every term of the
series~\eqref{s2_problems:eq:lap3d-exact} is harmonic and vanishes on the five faces, and
that the $b_n$ are the stated sine coefficients.
% src: research_benchmark/checks/verify_exact.out (heat1d 9.96e-07 with |u|max 0.590, wave1d 1.25e-04, wave2d 2.00e-07, laplace2d 4.44e-08, laplace3d 2.69e-07);
%      research_benchmark/scripts/verify_exact.py (2000 points drawn in the box shrunk by 5 per cent of each side at both ends; step h = 1e-4)
(ii)~Central finite-difference residuals of the coded solutions were evaluated at $2000$
random points of the inner 90 per cent of each box (every coordinate at least 5 per cent
of the side away from both ends), with step $h=10^{-4}$ in double precision. They are at
most $\sci{1.3}{-4}$ for \PWone, $\sci{1.0}{-6}$ for \PH\ and $\sci{2.7}{-7}$ for \PWtwo,
\PLtwo\ and \PLthree. The first two values are truncation errors of the difference
quotients, not errors of the formulas: for \PWone\ the mode of temporal frequency $8\pi$
leaves $\tfrac12(\omega^4-c^2k^4)h^2/12=\sci{1.25}{-4}$ with $\omega=8\pi$, $k=4\pi$,
$c^2=4$; for \PH\ the quotients leave $h^2(\pi^6/6+\pi^4/12)\abs{\uex}=\sci{1.68}{-6}\abs{\uex}$,
and $\abs{\uex}\le0.59$ at the sampled points, where $t\ge0.05$.
% (check by hand: the central first difference in t leaves (h^2/6) u_ttt = -(h^2/6) pi^6 u and the
%  second difference in x leaves (h^2/12) u_xxxx = (h^2/12) pi^4 u; 1.683e-6 * 0.590 = 9.93e-7,
%  against the reported maximum 9.96e-07.)
% src: research_benchmark/checks/verify_exact.out (laplace3d eval grid (64000 pts): max|series(4000)-series(8000)| = 0.00e+00)
(iii)~The series for \PLthree, truncated at $n=4000$ and at $n=8000$, gives identical
double-precision values at all $64\,000$ evaluation points.
% src: research_highdim/results/check_core.txt (laplace d=2,3,5,10: 0.00e+00; poisson d=2,3,5,10: at most 7.11e-15);
%      research_highdim/scripts/check_core.py (2000 points, float64)
(iv)~For \Pone\ and \Ptwo, automatic differentiation in double precision gives
$\max\abs{\lap\uex+f}\le\sci{7.11}{-15}$ over $2000$ random points for $d=2,3,5,10$, where $f$
is the right-hand side.
% src: data/closed_forms.json (floors.*.monte_carlo, 400 000 points; max difference 4.3e-4);
%      research_highdim/results/floors.json (50 000-point evaluation sets; max difference 1.8e-3)
(v)~The closed forms of Proposition~\ref{s2_problems:prop:floors} agree with Monte-Carlo
estimates on $400\,000$ points to within $\sci{5}{-4}$ and with the values on the evaluation
sets to within $\sci{2}{-3}$, for $d=2,3,5,10,20$.
